\chapter{حلّ المسائل}\label{ch:g3:problems}

معرفة العمليات الأربع شيء، ومعرفة \emph{أيّها} تحتاجه القصة شيء
آخر. يعطي هذا الفصل طريقة للمسائل اللفظية --- الخطوات الأربع
نفسها في كل مرة --- ويعلّم كيف نقرأ المعلومات من الجداول
والمخططات بالصور.

\section{الخطوات الأربع}

\begin{method}[حلّ مسألة لفظية]\label{met:g3:problems:steps}
\begin{enumerate}
  \item \emph{اقرأ} المسألة مرتين؛ وقل بكلماتك ما المطلوب،
        وما المعلوم.
  \item \emph{اختر} العملية: واُرسم خطاطة سريعة أو رسم أشرطة
        إن لم يكن الاختيار بيّنًا.
  \item \emph{احسب}، بالأعمدة أو ذهنيًّا، مع تقدير يكون شبكة
        أمان.
  \item \emph{أجب} بجملة كاملة، مع الوحدة --- وتحقّق من أن
        الجواب معقول (فالطفل لا يزن
        $300$ كغ).
\end{enumerate}
\end{method}

\begin{example}[أي عملية؟]\label{ex:g3:problems:which}
الكلمات تلمّح إلى العملية، والخطاطة تفصل:
\begin{itemize}
  \item \emph{الضمّ معًا / في المجموع}: جمع؛
  \item \emph{الأخذ / بكم يزيد / ما الناقص}:
        طرح؛
  \item \emph{كذا مجموعة من كذا / كذا مرة}:
        ضرب؛
  \item \emph{التقسيم بالعدل / كم مجموعة}: قسمة.
\end{itemize}
\end{example}

\begin{omfigure}
\begin{tikzpicture}[scale=0.62]
  % addition bar
  \draw[thick, fill=omDef!20] (0,0) rectangle (3.4,0.8);
  \draw[thick, fill=omThm!20] (3.4,0) rectangle (5.8,0.8);
  \node[font=\small] at (1.7,0.4) {$26$};
  \node[font=\small] at (4.6,0.4) {$17$};
  \draw[Stealth-Stealth] (0,-0.35) -- (5.8,-0.35);
  \node[below, font=\small] at (2.9,-0.4) {؟ في المجموع: $26 + 17$};
  % comparison bar
  \begin{scope}[xshift=9.4cm]
  \draw[thick, fill=omDef!20] (0,0.9) rectangle (5.2,1.7);
  \node[font=\small] at (2.6,1.3) {$52$};
  \draw[thick, fill=omThm!20] (0,0) rectangle (3.5,0.8);
  \node[font=\small] at (1.75,0.4) {$35$};
  \draw[Stealth-Stealth] (3.5,0.4) -- (5.2,0.4);
  \node[below, font=\small] at (2.6,-0.15)
    {؟ بكم يزيد: $52 - 35$};
  \end{scope}
\end{tikzpicture}

{\small رسوم الأشرطة تُظهر العملية: أجزاء متجاورة تدعو إلى
الجمع؛ وشريطان متقارَنان يدعوان إلى الطرح.}
\end{omfigure}

\begin{example}[مسألة من خطوتين]\label{ex:g3:problems:twostep}
«تشتري ميساء $3$ رزم من $6$ قنينات وتشرب $2$ منها. فكم
قنينة تبقى؟»
\begin{enumerate}
  \item المطلوب: عدد القنينات الباقية. والمعلوم: $3$
        رزم من $6$؛ و $2$ مشروبتان.
  \item الخطاطة: $3$ مجموعات من $6$، ثم أخذ $2$. عمليتان!
  \item الحساب: $3 \times 6 = 18$، ثم $18 - 2 = 16$.
  \item والجملة: تبقى $16$ قنينة.
\end{enumerate}
ولا تتوقف أبدًا بعد العملية الأولى: أعد قراءة السؤال لترى هل
انتهت القصة.
\end{example}

\begin{example}[عدد لا فائدة منه]\label{ex:g3:problems:trap}
«خبّاز عمره $52$ سنة يخبز $120$ خبزة ويبيع $85$.
فكم خبزة تبقى؟» العمر موضوع ليشتّت انتباهك: فالجواب لا يحتاج
إلا إلى $120 - 85 = 35$ خبزة. وقبل الحساب، فرّز الأعداد دائمًا:
أيّها يستعمله السؤال حقًّا؟
\end{example}

\section{قراءة الجداول والمخططات بالصور}

\begin{example}[جدول]\label{ex:g3:problems:table}
الكتب المستعارة من المكتبة:
\begin{center}
\renewcommand{\arraystretch}{1.3}
\begin{tabular}{l|cccc}
 & قصص مصورة & روايات & علوم & \omterm{def:g1:addition:def}{المجموع} \\
\hline
الأسبوع 1 & $14$ & $9$ & $6$ & $29$ \\
الأسبوع 2 & $11$ & $12$ & $8$ & $31$ \\
\end{tabular}
\end{center}
القراءة: الأسبوع 2، الروايات: $12$. والحساب من الجدول: في
الأسبوعين معًا استُعيرت $14 + 11 = 25$ قصة مصورة؛ وأكثر
الأسبوعين نشاطًا هو الأسبوع 2 ($31 > 29$).
\end{example}

\begin{example}[مخطط بالصور]\label{ex:g3:problems:pictogram}
في المخطط بالصور، يمثل كل رمز مقدارًا ثابتًا --- فاقرأ
المفتاح أولًا!

\begin{omfigure}
\begin{tikzpicture}[scale=0.62]
  \node[left, font=\small] at (0,2.4) {الاثنين};
  \foreach \i in {0,1,2} \fill[omDef] (0.5+\i*0.8,2.4) circle (0.26);
  \node[left, font=\small] at (0,1.2) {الثلاثاء};
  \foreach \i in {0,1,2,3,4} \fill[omDef] (0.5+\i*0.8,1.2) circle (0.26);
  \node[left, font=\small] at (0,0) {الأربعاء};
  \foreach \i in {0,1} \fill[omDef] (0.5+\i*0.8,0) circle (0.26);
  \fill[omDef] (0.5+1.6,0) circle (0.26);
  \node[right, font=\small] at (5.4,1.2)
    {المفتاح: نقطة واحدة $= 4$ كعكات مبيعة};
\end{tikzpicture}

{\small مخطط بالصور: الاثنين فيه $3$ نقط، إذن $3 \times 4 = 12$
كعكة؛ والثلاثاء $5$ نقط، أي $20$ كعكة؛ والأربعاء $3$ نقط، أي $12$
كعكة.}
\end{omfigure}
\end{example}

\section{تمارين}

\begin{exercise}[$\star$]\label{exo:g3:problems:1}
لكل قصة، \emph{سمِّ} العملية فقط (بلا حساب):
قطيع من $48$ خروفًا ينضم إلى قطيع من $37$؛ \ $56$ حبة كرز تُقسم
بين $7$ أطفال؛ \ $9$ علب من $8$ أقلام؛ \ كان عند علي $71$ بطاقة
فأعطى منها $18$.
\end{exercise}

\begin{exercise}[$\star$]\label{exo:g3:problems:2}
حُلّ بالخطوات الأربع: «في موقف $246$ سيارة في الصباح؛
وتصل $158$ أخرى ظهرًا. فكم سيارة فيه الآن؟»
\end{exercise}

\begin{exercise}[$\star$]\label{exo:g3:problems:3}
حُلّ: «حبل طوله $90$ م يُقطع إلى قطع من $9$ م. فكم
قطعة؟»
\end{exercise}

\begin{exercise}[$\star$]\label{exo:g3:problems:4}
حُلّ: «ثمن التذكرة $8$. فكم تكلّف $26$ تذكرة؟»
\end{exercise}

\begin{exercise}[$\star$]\label{exo:g3:problems:5}
اُرسم شريطًا ثم حُلّ: «عند هدى $64$ ملصقًا، وعند هشام أقلّ
بمقدار $29$. فكم عند هشام؟»
\end{exercise}

\begin{exercise}[$\star$]\label{exo:g3:problems:6}
خطوتان: «يسع الصندوق $45$ تفاحة. يستلم بقّال $6$ صناديق
ويبيع $178$ تفاحة. فكم تفاحة تبقى؟»
\end{exercise}

\begin{exercise}[$\star$]\label{exo:g3:problems:7}
بالاستعانة بجدول \cref{ex:g3:problems:table}: كم كتاب علوم
استُعير في الأسبوعين معًا؟ وبكم رواية يزيد الأسبوع 2 على
الأسبوع 1؟
\end{exercise}

\begin{exercise}[$\star$]\label{exo:g3:problems:8}
بالاستعانة بمخطط \cref{ex:g3:problems:pictogram}: كم كعكة
بيعت في الأيام الثلاثة معًا؟ وفي أي يوم بيع أكثر عدد من
الكعك؟
\end{exercise}

\begin{exercise}[$\star$]\label{exo:g3:problems:9}
في هذه المسألة عدد لا فائدة منه --- جِده ثم حُلّ: «كتاب من
$210$ صفحات ثمنه $12$. ويقرأ لؤي $35$ صفحة في اليوم. فكم يومًا
يلزمه لينهي الكتاب؟»
\end{exercise}

\begin{exercise}[$\star\star$]\label{exo:g3:problems:10}
«تزور المتحفَ أسرةٌ من $2$ من الكبار و $3$ من الأطفال. وثمن
تذكرة الكبير $9$، وتذكرة الطفل $4$. ويدفعون بورقة $50$.
فكم يستردّون؟» (ثلاث خطوات: الكبار، ثم الأطفال، ثم
\omterm{def:g3:division:remainder}{الباقي}.)
\end{exercise}

\begin{exercise}[$\star\star$]\label{exo:g3:problems:11}
اِبتكر مسألتك من خطوتين جوابها $100$، واكتب حلّها
بالخطوات الأربع، واختبرها على زميل.
\end{exercise}
